
\documentclass{article}
\usepackage{colortbl}
\usepackage{makecell}
\usepackage{multirow}
\usepackage{supertabular}

\begin{document}

\newcounter{utterance}

\centering \large Interaction Transcript for game `cladder', experiment `full\_v1.5\_default', episode 6252 with Qwen3.5{-}9B.
\vspace{24pt}

{ \footnotesize  \setcounter{utterance}{1}
\setlength{\tabcolsep}{0pt}
\begin{supertabular}{c@{$\;$}|p{.15\linewidth}@{}p{.15\linewidth}p{.15\linewidth}p{.15\linewidth}p{.15\linewidth}p{.15\linewidth}}
    \# & \multicolumn{2}{c}{Player} && \multicolumn{2}{c}{Game Master} \\
    \hline

    \theutterance \stepcounter{utterance}  
    & & & \multicolumn{4}{p{0.6\linewidth}}{
        \cellcolor[rgb]{0.9,0.9,0.9}{
            \makecell[{{p{\linewidth}}}]{
                \texttt{\tiny{[P1$\langle$GM]}}
                \texttt{You are an expert in causal inference. The following question is not a typical commonsense query, but rather a meticulously designed question created by a professor specializing in causal inference, intended to assess the students' mastery of the course content.} \\
\texttt{Imagine a self{-}contained, hypothetical world with only the following conditions, and without any unmentioned factors or causal relationships: Unobserved confounders has a direct effect on maternal smoking status and brown eyes. Cigarette tax has a direct effect on maternal smoking status. Maternal smoking status has a direct effect on brown eyes. Unobserved confounders is unobserved. We know that confounder active and high cigarette tax causes smoking mother. confounder active and smoking mother causes brown eyes. We observed the cigarette tax is high and confounder active. Would the infant has blue eyes if smoking mother instead of nonsmoking mother?} \\
\texttt{Start your answer with 'Yes' or 'No', followed by additional reasoning or evidence to support your explanation.} \\
            }
        }
    }
     \\ \\

    \theutterance \stepcounter{utterance}  
    & \multicolumn{4}{p{0.6\linewidth}}{
        \cellcolor[rgb]{0.561,0.737,0.561}{
            \makecell[{{p{\linewidth}}}]{
                \texttt{\tiny{[P1$\rangle$GM]}}
                \texttt{No} \\
\\ 
\texttt{**Reasoning:**} \\
\\ 
\texttt{In this hypothetical world, the causal structure is defined by the following relationships:} \\
\texttt{1.  **Unobserved Confounders (\$U\$)** \$\textbackslash{}rightarrow\$ **Maternal Smoking Status (\$S\$)**} \\
\texttt{2.  **Unobserved Confounders (\$U\$)** \$\textbackslash{}rightarrow\$ **Brown Eyes (\$E\$)**} \\
\texttt{3.  **Cigarette Tax (\$T\$)** \$\textbackslash{}rightarrow\$ **Maternal Smoking Status (\$S\$)**} \\
\texttt{4.  **Maternal Smoking Status (\$S\$)** \$\textbackslash{}rightarrow\$ **Brown Eyes (\$E\$)**} \\
\\ 
\texttt{The question asks for the counterfactual: Given that we observed \$T=\$ High and \$U=\$ Active (which resulted in \$S=\$ Smoking), would the infant have **blue** eyes if the mother had been **nonsmoking** (\$S=\$ Nonsmoking) instead?} \\
\\ 
\texttt{To answer this, we must determine the value of \$E\$ under the intervention \$do(S=\textbackslash{}text\{Nonsmoking\})\$ while holding the observed confounder \$U=\$ Active constant.} \\
\\ 
\texttt{1.  **Analyze the Direct Effect of Smoking on Eyes:** The prompt states that "Maternal smoking status has a direct effect on brown eyes." In causal diagrams, this implies a path \$S \textbackslash{}rightarrow E\$. However, the prompt *also* states that "Unobserved confounders has a direct effect on... brown eyes" (\$U \textbackslash{}rightarrow E\$).} \\
\texttt{2.  **Evaluate the Counterfactual Scenario:** We are intervening to set \$S = \textbackslash{}text\{Nonsmoking\}\$. This intervention blocks the causal path from smoking to eye color. However, the value of the eye color (\$E\$) is also determined by the unobserved confounder (\$U\$).} \\
\texttt{3.  **The Role of the Confounder (\$U\$):** We are given that \$U\$ is **Active**. The prompt explicitly states: "confounder active and smoking mother causes brown eyes." This phrasing suggests that the combination of \$U=\$ Active and \$S=\$ Smoking leads to the outcome "Brown Eyes." More critically, the structural equation for Eye Color includes a direct effect from \$U\$. Since \$U\$ is a direct cause of Eye Color and \$U\$ is fixed at "Active" in this scenario, the value of Eye Color is driven by \$U\$.} \\
\texttt{4.  **Determining the Outcome:** The prompt defines the outcome associated with "confounder active" as causing **brown eyes**. Even if we intervene to make the mother a nonsmoker (\$do(S=\textbackslash{}text\{Nonsmoking\})\$), the unobserved confounder \$U\$ remains **Active**. Because \$U\$ has a direct effect on eye color and "confounder active" is defined as a cause of brown eyes, the infant's eye color will remain **brown**. The intervention on smoking status does not alter the} \\
            }
        }
    }
    & & \\ \\

    \theutterance \stepcounter{utterance}  
    & & & \multicolumn{2}{p{0.3\linewidth}}{
        \cellcolor[rgb]{0.9,0.9,0.9}{
            \makecell[{{p{\linewidth}}}]{
                \texttt{\tiny{[GM$|$GM]}}
                \texttt{No} \\
            }
        }
    }
    & & \\ \\

    \theutterance \stepcounter{utterance}  
    & & & \multicolumn{2}{p{0.3\linewidth}}{
        \cellcolor[rgb]{0.9,0.9,0.9}{
            \makecell[{{p{\linewidth}}}]{
                \texttt{\tiny{[GM$|$GM]}}
                \texttt{no} \\
            }
        }
    }
    & & \\ \\

    \theutterance \stepcounter{utterance}  
    & & & \multicolumn{2}{p{0.3\linewidth}}{
        \cellcolor[rgb]{0.9,0.9,0.9}{
            \makecell[{{p{\linewidth}}}]{
                \texttt{\tiny{[GM$|$GM]}}
                \texttt{game\_result = WIN} \\
            }
        }
    }
    & & \\ \\

\end{supertabular}
}

\end{document}
